%% ----------------------------------------------------------------------------
%%  §6.1  Massimi e minimi relativi. Il teorema di Fermat
%% ----------------------------------------------------------------------------
\section{Massimi e minimi relativi. Il teorema di Fermat}
\label{sec:fermat}

Il teorema di Weierstrass (\cref{teo:weierstrass}) garantisce che una funzione continua
in un intervallo chiuso e limitato assume massimo e minimo, ma non dice \emph{dove}. In
questo paragrafo cominciamo a rispondere alla domanda. L'idea è semplice: in un punto
di massimo «interno» il grafico di una funzione derivabile non può né salire né
scendere, e quindi la retta tangente deve essere orizzontale. Per formulare l'idea in
modo preciso conviene introdurre una nozione di massimo più debole di quella usata fino
a ora, che guarda soltanto a ciò che accade \emph{vicino} al punto.

\begin{definizione}[Massimi e minimi relativi]\label{def:max-min-relativo}
\index{massimo!relativo}\index{minimo!relativo}\index{punto!di massimo relativo}
Sia $f:A\to\R$, con $A\subseteq\R$, e sia $x_{0}\in A$. Si dice che $x_{0}$ è un
\emph{punto di massimo relativo} (o \emph{locale}) per $f$ se esiste $\delta>0$ tale che
\begin{equation}\label{eq:max-relativo}
  f(x)\le f(x_{0}) \qquad\perogni x\in A\cap(x_{0}-\delta,x_{0}+\delta);
\end{equation}
si dice che $x_{0}$ è un \emph{punto di minimo relativo} (o \emph{locale}) per $f$ se
esiste $\delta>0$ tale che
\begin{equation}\label{eq:min-relativo}
  f(x)\ge f(x_{0}) \qquad\perogni x\in A\cap(x_{0}-\delta,x_{0}+\delta).
\end{equation}
Il valore $f(x_{0})$ si chiama, rispettivamente, \emph{massimo relativo} o \emph{minimo
relativo} di $f$. Se le disuguaglianze valgono in senso stretto per $x\neq x_{0}$, si
parla di massimo o minimo relativo \emph{stretto}. Se le disuguaglianze valgono per
\emph{ogni} $x\in A$ (non solo vicino a $x_{0}$), si parla di punto di massimo o di
minimo \emph{assoluto}. I punti di massimo e di minimo, relativo o assoluto, si dicono
anche \emph{punti di estremo}.
\end{definizione}

Ogni punto di massimo assoluto è anche di massimo relativo (basta prendere un $\delta$
qualunque), ma non viceversa. Nella \cref{fig:estremi-relativi} la funzione,
definita in $[a,b]$, ha un minimo relativo in $x_{1}$, che è anche il minimo assoluto,
un massimo relativo in $x_{2}$ e un minimo relativo in $x_{3}$, che non sono assoluti;
anche gli estremi dell'intervallo sono punti di estremo relativo: $a$ è di massimo
relativo (vicino ad $a$ i valori di $f$ sono minori di $f(a)$) e $b$ è il punto di
massimo assoluto. Si osservi che la definizione confronta $f(x_{0})$ soltanto con i
valori assunti \emph{nel dominio}: per questo ha senso anche negli estremi
dell'intervallo, dove l'intorno $(x_{0}-\delta,x_{0}+\delta)$ esce da $[a,b]$.

\begin{figure}[htb]
\centering
\begin{tikzpicture}
\begin{axis}[width=0.95\textwidth, height=5.6cm, xmin=0, xmax=6.8, ymin=0, ymax=3.8,
  axis lines=left, xtick={0.3,1.2,3.4,5.3,6.2},
  xticklabels={$a$,$x_{1}$,$x_{2}$,$x_{3}$,$b$}, ytick=\empty, clip=false]
  \fill[blu!8] (axis cs:2.6,0) rectangle (axis cs:4.2,3.55);
  \addplot[curva, domain=0.3:6.2, samples=200]
    {0.0875*x^4-1.155*x^3+4.9805*x^2-7.5684*x+4.72456};
  \addplot[tratteggio] coordinates {(2.6,2.8634) (4.2,2.8634)};
  \addplot[tratteggio] coordinates {(0.3,0) (0.3,2.8718)};
  \addplot[tratteggio] coordinates {(1.2,0) (1.2,1.0)};
  \addplot[tratteggio] coordinates {(3.4,0) (3.4,2.8634)};
  \addplot[tratteggio] coordinates {(5.3,0) (5.3,1.6031)};
  \addplot[tratteggio] coordinates {(6.2,0) (6.2,3.275)};
  \addplot[only marks, mark=*, mark size=1.6pt, rame] coordinates
    {(0.3,2.8718) (1.2,1.0) (3.4,2.8634) (5.3,1.6031) (6.2,3.275)};
  \node[font=\scriptsize, grigio, anchor=south] at (axis cs:3.4,3.56)
    {$(x_{2}-\delta,x_{2}+\delta)$};
  \node[font=\scriptsize, rame, above] at (axis cs:3.4,2.95) {max.\ rel.};
  \node[font=\scriptsize, rame, below right] at (axis cs:1.2,0.95) {min.\ ass.};
  \node[font=\scriptsize, rame, anchor=north west] at (axis cs:5.35,1.55) {min.\ rel.};
  \node[font=\scriptsize, rame, right] at (axis cs:6.25,3.275) {max.\ ass.};
  \node[font=\scriptsize, rame, anchor=south west] at (axis cs:0.3,2.95) {max.\ rel.};
\end{axis}
\end{tikzpicture}
\caption{Punti di estremo relativo e assoluto di una funzione definita in $[a,b]$. Nella
striscia evidenziata, intorno di $x_{2}$, il grafico sta sotto la retta orizzontale
$y=f(x_{2})$: $x_{2}$ è di massimo relativo, anche se la funzione assume altrove valori
maggiori di $f(x_{2})$. In $x_{1}$, $x_{2}$, $x_{3}$, interni all'intervallo, la tangente
è orizzontale; negli estremi $a$ e $b$ non lo è.}
\label{fig:estremi-relativi}
\end{figure}

La figura suggerisce la proprietà che stiamo cercando: nei punti di estremo
relativo \emph{interni} all'intervallo la tangente al grafico è orizzontale.

\begin{teorema}[di Fermat]\label{teo:fermat}\index{teorema!di Fermat}
Sia $f$ una funzione definita in un intervallo $[a,b]$ e sia $x_{0}\in(a,b)$ un punto
\emph{interno} di massimo o di minimo relativo per $f$. Se $f$ è derivabile in $x_{0}$,
allora
\begin{equation}\label{eq:fermat}
  f'(x_{0})=0 .
\end{equation}
\end{teorema}

\begin{schema}
\Ipotesi{\begin{itemize}[nosep, leftmargin=1.2em]
  \item $f:[a,b]\to\R$;
  \item $x_{0}\in(a,b)$ punto di massimo (o di minimo) relativo;
  \item $f$ derivabile in $x_{0}$.
\end{itemize}}
\Tesi{$f'(x_{0})=0$.}
\end{schema}

\begin{strategia}
Studiamo il \emph{segno} del rapporto incrementale $\frac{f(x_{0}+h)-f(x_{0})}{h}$ in un
punto di massimo. Il numeratore è $\le0$ per $h$ piccolo, qualunque sia il segno di $h$;
quindi il rapporto è $\le0$ se $h>0$ ed è $\ge0$ se $h<0$. Passando al limite, la
derivata destra è $\le0$ e la derivata sinistra è $\ge0$; ma la funzione è derivabile,
quindi le due derivate unilaterali coincidono con $f'(x_{0})$, che deve essere
contemporaneamente $\le0$ e $\ge0$. L'ipotesi che $x_{0}$ sia \emph{interno} serve a
poter prendere incrementi $h$ di entrambi i segni.
\end{strategia}

\begin{dimostrazione}
Supponiamo che $x_{0}$ sia un punto di massimo relativo (il caso del minimo si
riconduce a questo considerando la funzione $-f$, che ha un massimo relativo in $x_{0}$ e
derivata $-f'(x_{0})$). Per la \eqref{eq:max-relativo} esiste $\delta>0$ tale che
$f(x)\le f(x_{0})$ per ogni $x\in[a,b]\cap(x_{0}-\delta,x_{0}+\delta)$. Poiché $x_{0}$ è
interno, possiamo ridurre $\delta$ in modo che $\delta\le\min\{x_{0}-a,\,b-x_{0}\}$:
allora l'intero intervallo $(x_{0}-\delta,x_{0}+\delta)$ è contenuto in $[a,b]$, e
\begin{equation}\label{eq:fermat-numeratore}
  f(x_{0}+h)-f(x_{0})\le0 \qquad\perogni h\in(-\delta,\delta).
\end{equation}
Dividendo per $h$ si ottiene, per $0<h<\delta$ e per $-\delta<h<0$ rispettivamente,
\begin{equation}\label{eq:fermat-rapporti}
\begin{aligned}
  &\frac{f(x_{0}+h)-f(x_{0})}{h}\le0 \qquad\text{per } 0<h<\delta,\\[2pt]
  &\frac{f(x_{0}+h)-f(x_{0})}{h}\ge0 \qquad\text{per } -\delta<h<0.
\end{aligned}
\end{equation}
Poiché $f$ è derivabile in $x_{0}$, il rapporto incrementale tende a $f'(x_{0})$ sia per
$h\to0^{+}$ sia per $h\to0^{-}$ (\cref{prop:dx-sx}). Per il teorema del confronto dei
limiti (\cref{teo:confronto-funzioni}\,(ii)), il limite di una funzione $\le0$ è $\le0$,
e quello di una funzione $\ge0$ è $\ge0$; dalle \eqref{eq:fermat-rapporti} segue quindi
\begin{align*}
  f'(x_{0})&=\lim_{h\to0^{+}}\frac{f(x_{0}+h)-f(x_{0})}{h}\le0,\\
  f'(x_{0})&=\lim_{h\to0^{-}}\frac{f(x_{0}+h)-f(x_{0})}{h}\ge0 ,
\end{align*}
cioè $f'(x_{0})=0$.
\end{dimostrazione}

\begin{figure}[htb]
\centering
\begin{tikzpicture}
\begin{axis}[width=0.8\textwidth, height=5cm, xmin=-0.2, xmax=4.4, ymin=0.6, ymax=3.6,
  axis lines=left, xtick={1,2,3.2}, xticklabels={$x_{0}+h$,$x_{0}$,$x_{0}+h$},
  ytick=\empty, clip=false]
  \addplot[curva, domain=0:4, samples=100] {3-0.5*(x-2)^2};
  \addplot[curva3, domain=0.2:2.8] {3+0.5*(x-2)};
  \addplot[curva2, domain=1.2:4.2] {3-0.6*(x-2)};
  \addplot[tratteggio] coordinates {(0.4,3) (3.6,3)};
  \addplot[tratteggio] coordinates {(1,0.6) (1,2.5)};
  \addplot[tratteggio] coordinates {(3.2,0.6) (3.2,2.28)};
  \addplot[only marks, mark=*, mark size=1.6pt, blu] coordinates {(1,2.5) (2,3) (3.2,2.28)};
  \node[font=\scriptsize, verde!50!black, anchor=south east] at (axis cs:0.55,2.3) {$h<0$};
  \node[font=\scriptsize, rame, anchor=west] at (axis cs:4.25,1.68) {$h>0$};
\end{axis}
\end{tikzpicture}
\caption{La dimostrazione del teorema di Fermat. Vicino a un punto di massimo relativo
interno $x_{0}$, le secanti con $h<0$ (verde) hanno pendenza $\ge0$, quelle con $h>0$
(rame) pendenza $\le0$; al limite, la pendenza della tangente è insieme $\ge0$ e $\le0$.}
\label{fig:fermat}
\end{figure}

\begin{attenzione}[Il teorema di Fermat dà una condizione necessaria, non sufficiente]
Il viceversa del teorema di Fermat è falso: la derivata può annullarsi in un punto
che non è né di massimo né di minimo. L'esempio tipico è $f(x)=x^{3}$: risulta
$f'(0)=0$, ma $f(x)<0=f(0)$ per $x<0$ e $f(x)>0=f(0)$ per $x>0$, quindi in ogni intorno
di $0$ la funzione assume valori sia maggiori sia minori di $f(0)$. Inoltre il teorema
non dice nulla nei punti in cui $f$ non è derivabile: $f(x)=\abs x$ ha in $0$ il minimo
assoluto, ma non ha derivata in $0$ (\cref{es:modulo-non-derivabile}). Infine, l'ipotesi
che $x_{0}$ sia interno è essenziale: $f(x)=x$ in $[0,1]$ ha il massimo in $x_{0}=1$,
dove la derivata vale $1$.
\end{attenzione}

I punti $x_{0}$ in cui $f'(x_{0})=0$ si chiamano \emph{punti stazionari} (o
\emph{critici}) di $f$\index{punto!stazionario}. Che cosa accade negli estremi
dell'intervallo? Il ragionamento della dimostrazione si può ancora fare, ma con
incrementi di un solo segno, e fornisce una disuguaglianza invece di un'uguaglianza.

\begin{proposizione}[Estremi relativi e segno della derivata]\label{prop:fermat-estremi}
Sia $f:[a,b]\to\R$ e sia $x_{0}\in[a,b]$ un punto di massimo relativo per $f$, nel quale
$f$ è derivabile (con derivata unilaterale se $x_{0}=a$ oppure $x_{0}=b$). Allora
\begin{equation}\label{eq:fermat-max-generale}
  f'(x_{0})\,(x-x_{0})\le0 \qquad\perogni x\in[a,b].
\end{equation}
Se invece $x_{0}$ è un punto di minimo relativo, risulta
\begin{equation}\label{eq:fermat-min-generale}
  f'(x_{0})\,(x-x_{0})\ge0 \qquad\perogni x\in[a,b].
\end{equation}
In particolare: se $a$ è di massimo relativo allora $f'(a)\le0$, se $b$ è di massimo
relativo allora $f'(b)\ge0$; se $a$ è di minimo relativo allora $f'(a)\ge0$, se $b$ è di
minimo relativo allora $f'(b)\le0$; se $x_{0}\in(a,b)$ allora $f'(x_{0})=0$.
\end{proposizione}

\begin{schema}
\Ipotesi{$f:[a,b]\to\R$; \ $x_{0}\in[a,b]$ di massimo (minimo) relativo; \ $f$
derivabile in $x_{0}$ (unilateralmente se $x_{0}$ è un estremo).}
\Tesi{$f'(x_{0})(x-x_{0})\le0$ \ (rispettivamente $\ge0$) \ per ogni $x\in[a,b]$.}
\end{schema}

\begin{strategia}
Ripetiamo la dimostrazione del teorema di Fermat tenendo conto che, se $x_{0}=a$, sono
ammessi soltanto incrementi $h>0$, e se $x_{0}=b$ soltanto incrementi $h<0$. La forma
\eqref{eq:fermat-max-generale} riassume in una sola disuguaglianza tutti i casi: il
prodotto $f'(x_{0})(x-x_{0})$ è la variazione «prevista dalla tangente» quando ci si
sposta da $x_{0}$ a $x$, e in un punto di massimo non può essere positiva.
\end{strategia}

\begin{dimostrazione}
Sia $x_{0}$ di massimo relativo e sia $\delta>0$ come nella \eqref{eq:max-relativo}.
Se $x_{0}<b$, per $0<h<\min\{\delta,b-x_{0}\}$ il rapporto incrementale è $\le0$, e
passando al limite per $h\to0^{+}$ si ottiene $f'(x_{0})\le0$. Se $x_{0}>a$, per
$-\min\{\delta,x_{0}-a\}<h<0$ il rapporto incrementale è $\ge0$, e al limite per
$h\to0^{-}$ si ottiene $f'(x_{0})\ge0$. Quindi:
\begin{itemize}
  \item se $x_{0}=a$, vale solo il primo caso: $f'(a)\le0$, e poiché $x-a\ge0$ per ogni
  $x\in[a,b]$ risulta $f'(a)(x-a)\le0$;
  \item se $x_{0}=b$, vale solo il secondo caso: $f'(b)\ge0$, e poiché $x-b\le0$ si ha
  $f'(b)(x-b)\le0$;
  \item se $a<x_{0}<b$, valgono entrambi: $f'(x_{0})=0$, e la
  \eqref{eq:fermat-max-generale} è banalmente vera.
\end{itemize}
Il caso del minimo si ottiene applicando quanto provato alla funzione $-f$.
\end{dimostrazione}

\subsection{Un metodo per trovare massimi e minimi assoluti}

Combinando il teorema di Weierstrass con quello di Fermat otteniamo un procedimento
concreto. Sia $f$ continua in $[a,b]$ e derivabile in $(a,b)$, eccetto al più un numero
finito di punti. Per il teorema di Weierstrass $f$ ammette massimo e minimo assoluti;
un punto $x_{0}$ di massimo o di minimo assoluto è anche di estremo relativo, e
quindi, per il teorema di Fermat, deve appartenere a una delle tre categorie seguenti:
\begin{enumerate}[label=\textup{(\roman*)}]
  \item $x_{0}$ è uno degli estremi $a$, $b$ dell'intervallo;
  \item $x_{0}$ è un punto interno in cui $f$ non è derivabile;
  \item $x_{0}$ è un punto interno stazionario, cioè $f'(x_{0})=0$.
\end{enumerate}
Nei casi più comuni i punti delle tre categorie sono in numero finito: basta allora
calcolare i valori di $f$ in tali punti e confrontarli. Il più grande è il massimo
assoluto, il più piccolo è il minimo assoluto. Si noti che il ragionamento è
\emph{per esclusione}: non serve sapere a priori quali punti stazionari siano di
massimo o di minimo, perché l'esistenza degli estremi assoluti è garantita dal teorema
di Weierstrass. Il programma seguente esegue il procedimento con \texttt{sympy} per la
funzione $f(x)=x^{3}-3x$ nell'intervallo $[-2,3]$.

\begin{python}[Massimo e minimo assoluti in un intervallo chiuso]
import sympy as sp

x = sp.symbols("x", real=True)

def estremi_assoluti(f, a, b):
    """Massimo e minimo assoluti di f (derivabile) in [a, b]."""
    critici = sp.solveset(sp.diff(f, x), x, sp.Interval.open(a, b))
    candidati = sorted(set(critici) | {a, b})
    valori = {c: f.subs(x, c) for c in candidati}
    for c in candidati:
        print(f"  f({c}) = {valori[c]}")
    M, m = max(valori.values()), min(valori.values())
    dove = lambda y: [c for c in candidati if valori[c] == y]
    print("  max =", M, "in", dove(M))
    print("  min =", m, "in", dove(m))

print("f(x) = x^3 - 3x in [-2, 3]:")
estremi_assoluti(x**3 - 3*x, -2, 3)
\end{python}
\begin{risultato}
f(x) = x^3 - 3x in [-2, 3]:
  f(-2) = -2
  f(-1) = 2
  f(1) = -2
  f(3) = 18
  max = 18 in [3]
  min = -2 in [-2, 1]
\end{risultato}

Il massimo assoluto $18$ è assunto nell'estremo $b=3$, dove la derivata
$f'(3)=24$ non è nulla (ed è positiva, in accordo con la \cref{prop:fermat-estremi}); il
minimo assoluto $-2$ è assunto in due punti: nell'estremo $a=-2$, dove $f'(-2)=9\ge0$,
e nel punto stazionario interno $x=1$. Il punto stazionario $x=-1$ è di massimo
relativo, ma non assoluto. Ritroveremo il grafico di questa funzione nella
\cref{fig:cubica-monotonia}.

\begin{esempio}[La gittata massima]\label{es:gittata}
Un carico viene lanciato con velocità iniziale di modulo $v$, inclinata dell'angolo
$\theta$ sull'orizzontale. Trascurando la resistenza dell'aria, la distanza orizzontale
percorsa prima di tornare alla quota di partenza (la \emph{gittata}) è
\begin{equation}\label{eq:gittata}
  R(\theta)=\frac{v^{2}\sen2\theta}{g}, \qquad \theta\in\Bigl[0,\frac\pi2\Bigr],
\end{equation}
dove $g$ è l'accelerazione di gravità. La funzione $R$ è continua nell'intervallo
chiuso $[0,\frac\pi2]$, quindi ha massimo assoluto; è derivabile, con
$R'(\theta)=\frac{2v^{2}}{g}\cos2\theta$, che nell'intervallo aperto si annulla soltanto
per $2\theta=\frac\pi2$, cioè $\theta=\frac\pi4$. I candidati sono dunque $0$,
$\frac\pi4$, $\frac\pi2$, con $R(0)=R(\frac\pi2)=0$ e $R(\frac\pi4)=\frac{v^{2}}g$: la
gittata massima si ottiene con l'inclinazione di $45^{\circ}$, mentre i due estremi
sono punti di minimo assoluto. Coerentemente con la \cref{prop:fermat-estremi}, risulta
$R'(0)=\frac{2v^{2}}g\ge0$ (minimo nell'estremo sinistro) e
$R'(\frac\pi2)=-\frac{2v^{2}}g\le0$ (minimo nell'estremo destro).
\end{esempio}

\begin{python}[La gittata: punti critici e derivate negli estremi]
import sympy as sp

th, v, g = sp.symbols("theta v g", positive=True)
R = v**2 * sp.sin(2 * th) / g             # gittata nel vuoto
dR = sp.diff(R, th)
print("R'(theta) =", dR)
print("punti critici in (0, pi/2):",
      sp.solveset(dR, th, sp.Interval.open(0, sp.pi / 2)))
print("R'(0) =", dR.subs(th, 0),
      "   R'(pi/2) =", dR.subs(th, sp.pi / 2))
print("R(pi/4) =", R.subs(th, sp.pi / 4))
\end{python}
\begin{risultato}
R'(theta) = 2*v**2*cos(2*theta)/g
punti critici in (0, pi/2): {pi/4}
R'(0) = 2*v**2/g    R'(pi/2) = -2*v**2/g
R(pi/4) = v**2/g
\end{risultato}

\begin{idea}[L'efficienza massima, finalmente con tutti i crismi]
Nell'idea «Ritorno all'efficienza massima» del \cref{sec:operazioni-derivate} abbiamo
visto che la derivata dell'efficienza $E(C_{L})=\frac{C_{L}}{C_{D0}+kC_{L}^{2}}$ si annulla
in $C_{L}^{*}=\sqrt{C_{D0}/k}$. Ora possiamo concludere con il metodo appena descritto:
$E$ è continua in $[0,C_{L,\max}]$ (con $C_{L,\max}>C_{L}^{*}$), quindi ha massimo
assoluto; l'unico punto stazionario interno è $C_{L}^{*}$; negli estremi
$E(0)=0<E(C_{L}^{*})$ e $E(C_{L,\max})<E(C_{L}^{*})$, perché, come vedremo nel
\cref{sec:monotonia}, $E$ è strettamente decrescente per $C_{L}\ge C_{L}^{*}$. Il massimo assoluto è
quindi $E_{\max}=E(C_{L}^{*})=\frac1{2\sqrt{kC_{D0}}}$, in accordo con il
\cref{cap:funzioni-continue}.
\end{idea}

\begin{notastorica}[Fermat e il metodo dell'«adequalitas»]
Pierre de Fermat (1601?--1665), magistrato al parlamento di Tolosa e matematico per
passione, descrisse intorno al 1636--1638, in un breve scritto intitolato
\emph{Methodus ad disquirendam maximam et minimam} fatto circolare per lettera, un
procedimento per trovare i massimi e i minimi. In linguaggio moderno: si confronta
$f(x)$ con $f(x+E)$, si «adeguagliano» le due espressioni (\emph{adaequalitas}), si
divide per $E$ e infine si pone $E=0$. Nel primo esempio si divide un segmento di
lunghezza $b$ in due parti in modo che il rettangolo $x(b-x)$ sia massimo: da
$x(b-x)\sim(x+E)(b-x-E)$ segue $E(b-2x)-E^{2}\sim0$, poi $b-2x-E\sim0$, e per $E=0$ si
trova $x=\frac b2$: il rettangolo massimo è il quadrato. È, in sostanza, la condizione
$f'(x)=0$, ottenuta decenni prima del calcolo differenziale; ma dividere per $E$ e poi
porre $E=0$ non è passare al limite. Per i polinomi le due operazioni si equivalgono; se
invece $f$ contiene radicali, come $x+\sqrt{1-x}$, prima di dividere per $E$ bisogna
eliminarli elevando al quadrato, e nei casi meno semplici i calcoli diventano
proibitivi. Il nuovo calcolo superò l'ostacolo separando due momenti, il calcolo di $f'$
con regole generali e la soluzione di $f'(x)=0$: non a caso Leibniz vantava un metodo
«che non si arresta davanti alle quantità fratte o irrazionali» (epigrafe del
\cref{cap:derivate}). Già Keplero, nella \emph{Nova stereometria doliorum vinariorum}
(1615), aveva notato che vicino a un massimo le variazioni di una grandezza diventano
impercettibili. Fermat applicò lo stesso tipo di ragionamento all'ottica, enunciando nel
1662 il \emph{principio del tempo minimo}: la
luce, per andare da un punto a un altro, segue il percorso che richiede il tempo
minore (si veda l'\cref{es:snell}). Le sue opere furono pubblicate postume dal figlio
Samuel, nei \emph{Varia opera mathematica} (1679).
\end{notastorica}
